% TOP_PROBLEM: {"id":30007600,"problem_number":"LOCAL-30007600","title":"Climate risk versus green preferences","category_id":21,"category":{"id":21,"name":"economics","display_name":"Economics","description":"Open economic theory statements, published research agendas, and proposed scoped specifications; question provenance and historical priority are qualified per record.","slug":"economics","order_index":21,"created_at":"2026-10-08T00:00:00Z"},"status":"research_question","external_url":null,"legacy_ids":[],"published":false,"statement":"In the explicitly specified setting: How can researchers distinguish compensation for climate risk from investors’ green preferences when estimating green asset premia? The mathematical statement above fixes the target, assumptions and scope of this catalogue formulation.","economics":{"original_id":89,"field":"Asset pricing and financial markets","kind":"identification","formalization_basis":"adapted_research_question","source_status":"research_agenda","priority_status":"earliest_found","priority_note":"This is the earliest explicit statement verified in this audit; first priority for the full catalogue formulation is not established.","earliest_verified_reference":{"authors":["Bank of England"],"title":"Bank of England Agenda for Research, 2025–2028","year":2026,"url":"https://www.bankofengland.co.uk/research/bank-of-england-agenda-for-research","doi":null,"locator":"The Financial System / System-wide efficiency, final paragraph","evidence_summary":"Requests climate-risk quantification, pricing and disclosure research."},"current_reference":{"authors":["Klaus Adam","Stefan Nagel"],"title":"Expectations Data in Asset Pricing","year":2022,"url":"https://bpb-us-w2.wpmucdn.com/voices.uchicago.edu/dist/f/575/files/2022/01/ExpectationsSurvey_13.pdf","doi":"10.3386/w29977","locator":"Section 6, Future Research Directions, pp.35–36; abstract","evidence_summary":"Calls for belief aggregation and subjective-risk evidence."},"formalization_note":"The agenda asks climate-risk pricing questions; separating nonpecuniary preferences from risk uses a proposed intervention and is not a verbatim agenda definition."},"statusQualification":"Published question or research agenda verified at the cited locator. The mathematical specification is adapted for this catalogue and includes additional assumptions; source publication does not establish first-ever priority or certify this exact model remains unsolved. The agenda asks climate-risk pricing questions; separating nonpecuniary preferences from risk uses a proposed intervention and is not a verbatim agenda definition.","statusReviewedAt":"2026-10-08"}
% FIRST_POSTED: 2026-10-08
% ENTRY_KIND: statement_only
% !TeX program = lualatex
\documentclass[11pt]{article}
\usepackage{fontspec}
\usepackage[a4paper,margin=25mm]{geometry}
\usepackage{amsmath,amssymb,mathtools}
\usepackage{xurl}
\usepackage[hidelinks]{hyperref}
\newcommand{\sref}[2]{\hyperlink{source:#1:#2}{[S#2]}}
\setlength{\emergencystretch}{3em}
\begin{document}
% problemId:30007600
\section*{Climate risk versus green preferences}\label{30007600}
\subsection{Definitions and mathematical statement}
Fix green and otherwise comparable conventional securities with payoff vector \(X_{t+1}\), positive price vector \(P_t\) and investor holdings. Gross returns are \(R^k_{t+1}=X^k_{t+1}/P^k_t\) for security \(k\). Investor \(i\) has ordinary consumption utility and a nonpecuniary preference parameter \(g_i\) for financing environmentally beneficial activity; investors may also hold different beliefs about climate states. Let \(Z_{t+1}\) be a specified physical or transition climate-risk factor and \(M_{i,t+1}\) the consumption discount factor. Define a climate risk contribution using covariance of returns with \(M_{i,t+1}\) under the stated objective law; define a taste contribution as the equilibrium price change under \(do(g_i=0\text{ for all }i)\), keeping cash-flow technologies fixed. The task is to identify these distinct contributions to green valuation differences. Use stated assumptions linking surveys to beliefs and preferences, and variation that shifts tastes separately from climate exposures. Determine whether the joint law of prices, returns, holdings and survey responses distinguishes taste changes from anticipated climate cash-flow losses. If it does not, derive bounds or observational equivalence results. Account for firm responses when interpreting financing effects. A lower expected green return alone cannot identify either the climate-insurance motive or environmental preferences.

\par\textbf{Formulation and scope.} The agenda asks climate-risk pricing questions; separating nonpecuniary preferences from risk uses a proposed intervention and is not a verbatim agenda definition.

\par\textbf{Open-question provenance.} An explicit question or research agenda was verified at the source locator. This does not by itself certify that every later version remains unresolved \sref{30007600}{1}.

\par\textbf{Historical priority.} This is the earliest explicit statement verified in this audit; first priority for the full catalogue formulation is not established.

\subsection{Short English statement}
In the explicitly specified setting: How can researchers distinguish compensation for climate risk from investors’ green preferences when estimating green asset premia? The mathematical statement above fixes the target, assumptions and scope of this catalogue formulation.

\subsection{Sources}
\begin{itemize}\raggedright
\item[\textnormal{[S1]}]\hypertarget{source:30007600:1}{} Bank of England. 2026. Bank of England Agenda for Research, 2025–2028. The Financial System / System-wide efficiency, final paragraph.
\url{https://www.bankofengland.co.uk/research/bank-of-england-agenda-for-research}.
{\small\textit{Earliest verified question/agenda reference; historical priority qualified below. Requests climate-risk quantification, pricing and disclosure research.}}

\item[\textnormal{[S2]}]\hypertarget{source:30007600:2}{} Klaus Adam, Stefan Nagel. 2022. Expectations Data in Asset Pricing. Section 6, Future Research Directions, pp.35–36; abstract.
\url{https://bpb-us-w2.wpmucdn.com/voices.uchicago.edu/dist/f/575/files/2022/01/ExpectationsSurvey_13.pdf}.
DOI: 10.3386/w29977.
{\small\textit{Supporting or current literature; see evidence qualification. Calls for belief aggregation and subjective-risk evidence.}}
\end{itemize}
\end{document}
